In order to establish the premise for the interaction, this subsection will discuss several key terms related to particle collision kinematics. Firstly, we shall define $-Q^2$ as the square of the momentum transfer between incoming particle and the target particle. In case of a photon with four-momentum $q^\mu$, one would get $Q^2 = -q^2$. Another important quantity is Bjorken x , or x for short, which is defined as:

\begin{equation}\label{eq:bjorken_x}
	x = \frac{Q^2}{2pq}\,,
\end{equation}

where $p$ is the momentum of the target. Bjorken x can be thought of as the fraction of momentum carried by the quark which interacted with the photon.